\section{Théorème de Gauss, Identité de Bézout, Algorithme d’Euclide}

\subsection{Motivation}
La théorie des nombres ne se contente pas d'étudier les propriétés abstraites des entiers ; elle fournit les algorithmes qui sécurisent le monde numérique. L'algorithme d'Euclide est l'un des plus vieux algorithmes au monde, et sa version "étendue" (liée à l'identité de Bézout) permet de calculer des inverses modulaires. Ces calculs sont le cœur mathématique du chiffrement RSA, utilisé quotidiennement pour sécuriser les transactions bancaires et les communications sur Internet.

\subsection{Approche par problème}
Imaginez que vous disposiez d'un récipient de 5 litres et d'un autre de 3 litres, sans aucune graduation. Comment faire pour obtenir exactement 1 litre d'eau ? 

Mathématiquement, cela revient à chercher deux entiers $u$ et $v$ (le nombre de fois où l'on remplit ou vide chaque récipient) tels que $5u + 3v = 1$. Cette équation, appelée équation diophantienne linéaire, possède des solutions entières uniquement grâce aux relations profondes qui lient le nombre 5, le nombre 3, et leur PGCD. Le théorème de Bézout et l'algorithme d'Euclide nous donnent la méthode systématique pour trouver ces solutions.

\subsection{Définitions et théorèmes}

\subsubsection{L'Algorithme d'Euclide}

L'algorithme d'Euclide permet de calculer le PGCD de deux entiers sans avoir à chercher tous leurs diviseurs. Il repose sur le lemme suivant : si $a = bq + r$, alors $\text{PGCD}(a, b) = \text{PGCD}(b, r)$.

\begin{theoreme}[Algorithme d'Euclide]
Soient $a$ et $b$ deux entiers naturels non nuls avec $a > b$. En effectuant des divisions euclidiennes successives :
\begin{align*}
a &= bq_1 + r_1 \quad (0 < r_1 < b) \\
b &= r_1q_2 + r_2 \quad (0 < r_2 < r_1) \\
r_1 &= r_2q_3 + r_3 \quad (0 < r_3 < r_2) \\
&\dots \\
r_{n-2} &= r_{n-1}q_n + r_n \quad (0 < r_n < r_{n-1}) \\
r_{n-1} &= r_nq_{n+1} + 0
\end{align*}
Le dernier reste non nul, $r_n$, est exactement le $\text{PGCD}(a, b)$.
\end{theoreme}



[Image of flowchart of the Euclidean algorithm]


\subsubsection{Identité de Bézout}

\begin{theoreme}[Identité de Bézout]
Deux entiers relatifs $a$ et $b$ sont premiers entre eux si et seulement s'il existe deux entiers relatifs $u$ et $v$ tels que :
$$au + bv = 1$$
\end{theoreme}

\begin{propriete}[Généralisation - Théorème de Bachet-Bézout]
Soient $a$ et $b$ deux entiers non nuls. Si $d = \text{PGCD}(a, b)$, alors il existe deux entiers relatifs $u$ et $v$ tels que :
$$au + bv = d$$
De plus, l'équation $ax + by = c$ admet des solutions entières $(x, y)$ si et seulement si $c$ est un multiple de $d$.
\end{propriete}

\subsubsection{Théorème de Gauss}

\begin{theoreme}[Théorème de Gauss]
Soient $a$, $b$ et $c$ trois entiers relatifs non nuls. 
Si $a$ divise le produit $bc$ et si $a$ et $b$ sont premiers entre eux ($\text{PGCD}(a, b) = 1$), alors $a$ divise $c$.
\end{theoreme}

\subsection{Exemples de difficulté croissante}

\begin{exemple}[Recherche de PGCD par l'algorithme d'Euclide]
Déterminer le PGCD de $420$ et $135$.
\end{exemple}

\textbf{Correction :} On applique l'algorithme d'Euclide par divisions successives :
\begin{itemize}
    \item $420 = 135 \times 3 + 15$ (le reste est 15)
    \item $135 = 15 \times 9 + 0$ (le reste est 0)
\end{itemize}
Le dernier reste non nul est $15$. Donc, $\text{PGCD}(420, 135) = 15$.

\begin{exemple}[Résolution d'une équation avec Bézout et Gauss]
Résoudre dans $\mathbb{Z}^2$ l'équation $(E) : 5x - 3y = 2$.
\end{exemple}

\textbf{Correction :} 
1. \textbf{Recherche d'une solution particulière :} On remarque de manière évidente que $5 \times 1 - 3 \times 1 = 2$. Donc le couple $(1, 1)$ est une solution particulière.
2. \textbf{Soustraction :} Soit $(x, y)$ une solution. On a :
$$5x - 3y = 2$$
$$5(1) - 3(1) = 2$$
En soustrayant ligne à ligne : $5(x - 1) - 3(y - 1) = 0$, soit $5(x - 1) = 3(y - 1)$.
3. \textbf{Application du théorème de Gauss :} $5$ divise $3(y - 1)$. Or $\text{PGCD}(5, 3) = 1$. D'après le théorème de Gauss, $5$ divise $(y - 1)$.
Il existe donc un entier $k$ tel que $y - 1 = 5k$, d'où $y = 5k + 1$.
4. \textbf{Substitution :} On remplace $y - 1$ par $5k$ dans l'égalité $5(x - 1) = 3(y - 1)$ :
$$5(x - 1) = 3(5k) \implies x - 1 = 3k \implies x = 3k + 1$$
Les solutions sont donc les couples de la forme $(3k + 1, 5k + 1)$ avec $k \in \mathbb{Z}$.

\subsection{Exercice pratique avec Python}

\textbf{Objectif :} Implémenter l'algorithme d'Euclide étendu. 

Cet algorithme ne se contente pas de trouver le PGCD, il "remonte" les calculs pour trouver les coefficients $u$ et $v$ de l'identité de Bézout ($au + bv = \text{PGCD}$). C'est indispensable pour calculer des clés cryptographiques.

\begin{verbatim}
def euclide_etendu(a, b):
    """
    Retourne (pgcd, u, v) tels que a*u + b*v = pgcd
    basé sur l'algorithme d'Euclide étendu.
    """
    if b == 0:
        return a, 1, 0
    else:
        # Appel récursif
        d, u_prime, v_prime = euclide_etendu(b, a % b)
        
        # On reconstitue les coefficients u et v
        u = v_prime
        v = u_prime - (a // b) * v_prime
        
        return d, u, v

# --- Tests ---
a, b = 420, 135
pgcd, u, v = euclide_etendu(a, b)

print(f"Pour a={a} et b={b} :")
print(f"PGCD trouvé : {pgcd}")
print(f"Coefficients de Bézout : u={u}, v={v}")
print(f"Vérification : {a}*({u}) + {b}*({v}) = {a*u + b*v}")

# Test avec le problème d'introduction (récipients de 5L et 3L)
d, u, v = euclide_etendu(5, 3)
print("\nProblème des récipients (5L et 3L) :")
print(f"5*({u}) + 3*({v}) = {d}")
# Interprétation : Remplir 2 fois le récipient de 5L (10L), 
# et vider 3 fois son contenu dans celui de 3L (-9L) laisse 1L.
\end{verbatim}